\documentclass[10pt]{article}
\usepackage[accepted]{icml2026}

\usepackage[T1]{fontenc}
\usepackage[utf8]{inputenc}
\usepackage{lmodern}
\usepackage{microtype}
\usepackage{amsmath,amssymb}
\usepackage{booktabs}
\usepackage{array}
\usepackage{enumitem}
\usepackage[numbers,sort&compress]{natbib}
\usepackage{xcolor}
\usepackage{hyperref}
\usepackage{float}

\hypersetup{
  colorlinks=true,
  linkcolor=blue!50!black,
  citecolor=blue!50!black,
  urlcolor=blue!50!black
}

\newcommand{\RESULTPENDING}{\textcolor{red!70!black}{\textbf{RESULT PENDING}}}

\icmltitlerunning{Tail-Aware Prediction Market Forecasting}

\title{\textbf{Calibration Is Necessary but Not Sufficient: Tail-Aware Prediction Market Forecasting}}
\author{}
\date{}

\begin{document}

\twocolumn[
\maketitle
\begin{center}
\vspace{-1.75em}
\begin{minipage}{0.92\textwidth}
\begin{abstract}
Prediction markets are often evaluated descriptively: a market is called
well-calibrated when prices match realized frequencies on average. That view is
useful but incomplete. Calibration is not a primitive property of a market; it
is an outcome generated by liquidity, participation, timing, category
structure, and the way information arrives. We propose a causal study of those
drivers using a continuously collected Polymarket research stack. The local
repository already supports market ingestion, trade aggregation, calibration
tracking, market microstructure computation, and several causal-analysis
modules. Our empirical strategy centers on resolved binary markets, deriving
pre-resolution probabilities from trade histories and estimating how calibration
varies with market structure, activity, and domain. The paper's central claim
is that prediction-market calibration should be modeled as a causal outcome with
heterogeneous mechanisms rather than as a single pooled descriptive statistic.
In a first broad pass with a $\geq 10$ trade filter, the current pipeline yields
51{,}023 resolved binary markets, with final observed pre-resolution prices
available for 45{,}506 of them. On that slice, final-price calibration is
already strong in aggregate (Brier $= 0.0690$, ECE$_{10}=0.0196$), but the
descriptive heterogeneity across categories is large enough to motivate a causal
follow-up. We further argue that binary calibration is necessary but not
sufficient in fat-tailed or asymmetric domains: when multiple threshold markets
exist on the same event, the relevant object is the exceedance curve rather than
a single threshold probability. A new ladder study on 1{,}144 distinct
threshold ladders shows that this extension is already empirically available in
the repo: 16.1\% of ladders exhibit some adjacent-strike dominance violation
and 5.2\% exhibit a material ($>$2c) violation, while at-the-money contract
accuracy has near-zero correlation with adjacent arbitrage edge ($r=0.083$).
The current draft is therefore grounded in implemented infrastructure and live
empirics rather than in a purely speculative proposal.
\end{abstract}
\end{minipage}
\end{center}
\vspace{0.6em}
]

\section{Introduction}

Prediction markets are increasingly treated as live probabilistic forecasts for
politics, macroeconomics, sports, crypto, and science. Classical work shows
that these markets can be reasonably calibrated, especially near resolution, but
it also shows that calibration varies with horizon and market design
\citep{page2013prediction}. More recent work suggests that calibration is
structured and domain-specific rather than uniform across all markets
\citep{le2026crowd}. What remains underdeveloped is a causal account of why some
markets calibrate well while others fail.

That gap matters for both market science and AI forecasting. If a forecaster
uses market prices as a baseline, it needs to know when the market should be
trusted, when it should be discounted, and which observed market conditions are
mechanisms rather than symptoms. A purely correlational treatment of spread,
volume, or concentration is not enough. For actual trading systems, the target
is even stricter: not just forecast quality, but deployable edge under Kelly
sizing, portfolio limits, and tail-risk controls.

There is also a deeper objection. In thin-tailed settings, good calibration of a
single probability forecast may already be close to the decision-relevant
object. In fat-tailed or asymmetric settings, however, a well-calibrated
threshold probability can still be a poor guide to expected loss, tail
exposure, or right-tail opportunity. Taleb et al.\ argue that tournament-style
probability assessment and tail-risk assessment can come apart sharply
\citep{taleb2023probability}. For nonnegative outcomes with exceedance curve
$S(k)=P(X>k)$, the decision-relevant objects are integrals of that curve such
as $E[X]=\int S(k)\,dk$ and $E[X^2]=2\int kS(k)\,dk$, not merely one point on
$S$. That critique is especially relevant when the forecast is later used
inside a trading, hedging, or policy system.

This paper targets that gap. Our central question is:
\begin{quote}
Which market-level and microstructural features causally drive
prediction-market calibration, and how does that structure vary across
domains?
\end{quote}
The present draft adds a second question that follows directly from the first:
\begin{quote}
When the underlying outcome is naturally thresholded at many levels, can we
recover a useful exceedance curve rather than scoring only one binary contract?
\end{quote}

\paragraph{Contributions.}
\begin{enumerate}[leftmargin=1.5em,itemsep=0.2em]
  \item We reframe prediction-market calibration as an outcome generated by
  market structure and microstructure rather than as a single descriptive
  statistic.
  \item We define a repo-grounded extraction pipeline for resolved Polymarket
  markets using trade histories and metadata instead of the currently unreliable
  historical price table.
  \item We specify a causal workflow that combines calibration slicing,
  structure learning, and effect estimation.
  \item We operationalize ladder-level exceedance diagnostics, including
  truncated first- and second-moment proxies and a variance proxy recovered
  from Over/Under strike ladders.
  \item We connect ladder incoherence to trading deployment: cross-strike
  dominance violations imply relative-value spreads and supply a natural
  risk-management label for Kelly haircuts or portfolio gating.
\end{enumerate}

\section{Repo-Grounded Data and Measurement}

The drafting strategy is intentionally constrained by what is already
implemented in the local \texttt{polymarket} repository. The codebase contains a
continuous Polymarket ingestion pipeline, calibration primitives,
microstructure-computation jobs, and a family of causal-analysis modules that
already implement PC-style structure learning, transfer-entropy analysis,
interrupted time-series event-impact analysis, and manipulation detection.

The strongest current empirical substrate is the resolved binary market set. For
each market, we define the \emph{positive} event as the first listed outcome.
Each trade is then mapped into a positive-event probability:
\[
\hat p_t =
\begin{cases}
\texttt{price}_t & \text{if trade outcome} = \text{positive outcome} \\
1-\texttt{price}_t & \text{otherwise.}
\end{cases}
\]
This lets us reconstruct a probability path even when only one side of the
binary market is actively traded.

Our contract-level outcome is Brier error at one or more horizons before
resolution:
\[
\text{Brier}_h = (\hat p_h - y)^2,
\]
where $\hat p_h$ is the last observed implied probability before horizon $h$ and
$y \in \{0,1\}$ indicates whether the positive outcome ultimately wins
\citep{brier1950verification}. We also compute log loss, Expected Calibration
Error, and binned reliability curves.

\paragraph{Why trade-derived probabilities?}
The checked-in audit in the repo flags the historical \texttt{market\_prices}
table as unreliable for retrospective analysis because \texttt{condition\_id}
was empty in that snapshot. The first empirical pass therefore reconstructs
pre-resolution prices from \texttt{market\_trades}, which do contain
\texttt{condition\_id}, \texttt{outcome}, \texttt{price}, and
\texttt{timestamp}.

\paragraph{Current extraction footprint.}
On April 14, 2026, the first calibration pass with a $\geq 10$ trade filter
returned 51{,}023 resolved binary markets across 343 categories. Coverage is
45{,}506 markets at the final observed pre-resolution price, 35{,}488 at 6
hours before resolution, 25{,}231 at 24 hours, and 6{,}585 at 168 hours.
Median trade count is 59 and median market duration is 4.30 days.

\paragraph{Beyond single-threshold calibration.}
A binary contract gives one exceedance probability, $P(X > k)$. For thin-tailed
questions this may already be useful, but for fat-tailed or strongly asymmetric
domains it is not enough: the decision-relevant object is closer to an
exceedance curve, a tail integral, or a scenario-conditional distribution than
to one threshold probability. In practical terms, this means Proposal A should
not stop at ``is this contract calibrated?'' but should also ask whether the
market provides enough threshold information to approximate more of the right
tail.

When a ladder of thresholds $k_1 < \cdots < k_n$ is available, the exceedance
curve can be converted into truncated moment proxies:
\[
\widehat{E[X]} \approx k_1 + \sum_i S(k_i)(k_{i+1}-k_i), \qquad
\widehat{E[X^2]} \approx k_1^2 + \sum_i S(k_i)(k_{i+1}^2-k_i^2),
\]
which in turn give a truncated variance proxy
$\widehat{\mathrm{Var}}(X)=\widehat{E[X^2]}-\widehat{E[X]}^2$. A single
threshold contract cannot recover these objects on its own.
The same ladders can also be checked for dominance violations: if
$P(X>k_{i+1}) > P(X>k_i)$ for $k_{i+1} > k_i$, then the ladder is not merely
poorly calibrated but internally inconsistent in a way that matters for
relative-value trading and portfolio risk control.

\section{Feature Families}

The paper separates broadly available market features from sparse high-frequency
features.

\paragraph{Primary feature families.}
These are available for most resolved markets with sufficient trade data:
\begin{itemize}[leftmargin=1.5em,itemsep=0.2em]
  \item market structure: total volume, duration, outcome count, negative-risk
  structure, category;
  \item trade activity: trade count, average trade size, maximum trade size,
  trade concentration, trades per day;
  \item temporal concentration: late-volume ratio and recent activity proxies;
  \item recent market movement: one-day and one-week price-change fields from
  market metadata.
\end{itemize}

\paragraph{Secondary feature families.}
These exist in the repo but are currently much sparser on resolved markets:
\begin{itemize}[leftmargin=1.5em,itemsep=0.2em]
  \item bid-ask spread, effective spread, realized spread;
  \item order-book imbalance and depth metrics;
  \item wallet concentration and insider-score features;
  \item toxicity, price impact, and Kyle's lambda.
\end{itemize}
These features are scientifically important, but the current paper treats them
as a subset analysis unless historical coverage is improved before submission.

\section{Identification Strategy}

The workshop memo proposed a four-stage pipeline. The current paper narrows that
plan to what can be defended from the code and data already in hand.

\paragraph{Stage 1: descriptive calibration decomposition.}
We first quantify how calibration changes with category, trade intensity, volume
regime, and horizon. This yields the first hard empirical result and motivates
the causal question with visible heterogeneity.

\paragraph{Stage 2: structure learning.}
The repo already contains simplified PC-style graph discovery and complementary
directional information-flow modules. The first submission therefore estimates a
feature-level graph over observed market attributes and calibration outcomes,
with category-specific re-estimation where sample sizes permit
\citep{spirtes2000causation,pearl2009causality}.

\paragraph{Stage 3: effect estimation.}
After selecting interpretable treatments with reasonable overlap, we estimate
their effect on calibration. The most realistic first-pass treatments are:
high versus low trade intensity, high versus low concentration, short versus
long time-to-resolution, and high versus low recent information-arrival
regimes.

\paragraph{Stage 4: validation.}
Validation is practical rather than maximalist: temporal holdout by resolution
date, category-specific robustness checks, and analyses with and without sparse
high-frequency features.

\paragraph{Distributional extension via ladders.}
When Polymarket lists many Over/Under thresholds on the same event, the market
is implicitly offering multiple samples from a survival curve. This creates a
second empirical layer on top of contract-level calibration:
\begin{itemize}[leftmargin=1.5em,itemsep=0.2em]
  \item monotonicity diagnostics for the implied exceedance curve;
  \item integrated Brier score across thresholds;
  \item truncated first- and second-moment proxies, plus a variance proxy;
  \item decoupling tests between at-the-money contract accuracy and full-ladder quality;
  \item adjacent-strike arbitrage labels for profitability and risk gating;
  \item eventual tail-shape tests on the right tail where enough strikes exist.
\end{itemize}
This is not yet full extreme-value theory, but it is a concrete way to move
beyond probability-at-one-threshold and toward the object that matters for
expected-value and tail-risk decisions.

\section{Initial Empirical Status}

The empirical work has started. The new script
\texttt{experiments/run\_calibration\_study.py} extracts resolved binary markets
from ClickHouse, reconstructs positive-outcome probability paths from trades,
and writes paper-facing artifacts: market-level CSVs, horizon-wise summary
metrics, reliability bins, and slice tables by category, trade count, and
volume quartile.

The first completed pass yields the following overall calibration summaries:
\begin{itemize}[leftmargin=1.5em,itemsep=0.2em]
  \item final observed price: $n=45{,}506$, Brier $=0.0690$,
  log loss $=0.2080$, ECE$_{10}=0.0196$;
  \item 6h horizon: $n=35{,}488$, Brier $=0.0896$,
  ECE$_{10}=0.0177$;
  \item 24h horizon: $n=25{,}231$, Brier $=0.0549$,
  ECE$_{10}=0.0193$;
  \item 168h horizon: $n=6{,}585$, Brier $=0.0799$,
  ECE$_{10}=0.0095$.
\end{itemize}
These horizon summaries are descriptive and are not directly comparable as
time-to-resolution effects because each horizon is observed on a different
coverage subset.

The key descriptive finding is heterogeneity. At the final-price horizon,
Esports markets have Brier $=0.0230$ on $8{,}763$ contracts, while Sports
markets have Brier $=0.1778$ on $12{,}999$ contracts; Politics sits between
those extremes at Brier $=0.0577$ on 844 contracts. Trade-count bins are much
more stable, clustering between roughly 0.067 and 0.071 Brier, while the top
volume quartile is noticeably worse than the bottom quartile (0.0855 vs.\ 0.0516).

These are still descriptive rather than causal results, but they already show
that calibration is not homogeneous across the Polymarket universe. The code
path now exists in the repo and runs directly on the live ClickHouse data
backing the rest of the stack, which means the next step is no longer dataset
construction but effect estimation and graph recovery.

\paragraph{Threshold-ladder extension.}
To address the limitation of single-contract calibration, we added a second
script, \texttt{experiments/run\_threshold\_ladder\_study.py}, that reuses the
market-level calibration output and reconstructs Over/Under ladders from
multiple thresholds on the same event. To avoid mixing unrelated props inside
one event slug, ladders are grouped by normalized question template. On the
current data snapshot, this returns 13{,}166 O/U markets and 1{,}144 distinct
ladders with at least three thresholds. The median ladder has 4 thresholds and
the mean has 4.33.

At the ladder level, the mean integrated final-price Brier score is 0.1843.
The mean isotonic L1 repair gap is 0.0030, 16.1\% of ladders exhibit at least
one adjacent dominance violation, and 5.2\% exhibit a material ($>$2c)
violation. Most ladders remain too short to distinguish cleanly between
exponential and power-law right-tail decay, but the infrastructure now exists
to study exactly that question when richer strike ladders appear.

The stronger empirical point is deployability. The updated ladder script now
extracts a central at-the-money contract for each ladder together with
cross-strike arbitrage labels. Mean at-the-money Brier is 0.2004. It correlates
strongly enough with full-ladder Brier to remain useful ($r=0.679$), but it has
almost no relationship with adjacent arbitrage gross edge ($r=0.083$). Even
among bottom-quartile at-the-money ladders, 31.5\% still contain some adjacent
dominance violation and 6.6\% still contain a material ($>$2c) one. For
trading systems, that means calibration alone is not the right deployment
label: Kelly sizing and portfolio limits should also react to ladder coherence.

\section{Limitations}

The first submission must stay honest about data constraints.
\begin{enumerate}[leftmargin=1.5em,itemsep=0.2em]
  \item Historical price reconstruction currently depends on trades rather than
  a clean dedicated price table.
  \item Orderbook overlap for resolved markets is much thinner than trade
  overlap, so orderbook-heavy claims should remain subset-specific.
  \item Wallet-derived participation features are missing for a substantial
  fraction of markets.
  \item Category imbalance is severe; pooled estimates can be dominated by
  sports and esports unless explicitly stratified.
  \item The threshold-ladder analysis is currently dominated by sports and
  esports totals, which are useful for methodology but not yet a direct test of
  catastrophic fat-tail forecasting.
  \item The recovered moment and variance proxies are truncated to the observed
  strike range and are not yet normalized across categories with different
  outcome units.
  \item Some large ladder violations likely reflect stale last-trade states in
  thin contracts, so observed relative-value spreads are a profitability label
  rather than a guarantee of executable arbitrage at scale.
\end{enumerate}

\section{Conclusion}

Prediction-market calibration should not be treated as a single descriptive
statistic detached from the mechanisms that generate it. The local Polymarket
stack now makes it possible to study calibration as a market-level causal
outcome, beginning with resolved binary markets and trade-derived implied
probabilities. The immediate next step is to move from these descriptive
calibration slices to effect estimation and graph recovery on the most plausible
drivers of market quality. At the same time, the threshold-ladder results push
Proposal A beyond plain calibration: when the underlying object of interest is
tail-sensitive, the right target is not a single well-calibrated threshold
probability but a better estimate of the exceedance curve or its tail
functionals. For Polymarket trading systems, the same lesson appears in more
operational form: deployment should depend on ladder coherence and
distributional quality, not only on one-contract calibration. A natural next
step is to combine these diagnostics with scenario-conditioned forecasting
questions of the kind used in recent AI economics worldbuilding exercises
\citep{convergence2024threshold2030}.

\bibliographystyle{plainnat}
\bibliography{references}

\end{document}
